\section{Technical Features}\label{sec:technical-features}

\subsection{Code Listings}\label{subsec:code-listings}
The \texttt{listings} package provides excellent support for displaying code with syntax highlighting.

\subsubsection{Basic Code Listing}\label{subsubsec:basic-listing}
To include code in your document, use the \texttt{lstlisting} environment:

\begin{lstlisting}[language=Python, caption=Basic Python function,label={lst:lstlisting}]
def factorial(n):
    """Calculate factorial of n."""
    if n <= 1:
        return 1
    else:
        return n * factorial(n-1)

# Calculate factorial of 5
result = factorial(5)
print(f"5! = {result}")  # Output: 5! = 120
\end{lstlisting}

\subsubsection{Customizing Listings}\label{subsubsec:customized-listing}
You can customize the appearance of your code listings:

\begin{lstlisting}[
    language=C++,
    caption=C++ Class Example,
    basicstyle=\small\ttfamily,
    keywordstyle=\color{blue},
    commentstyle=\color{green!60!black},
    stringstyle=\color{red},
    numbers=left,
    numberstyle=\tiny,
    breaklines=true,
    frame=single
]
#include <iostream>
using namespace std;

class Rectangle {
private:
    double width;
    double height;

public:
    // Constructor
    Rectangle(double w, double h) {
        width = w;
        height = h;
    }

    // Calculate area
    double getArea() {
        return width * height;
    }

    // Calculate perimeter
    double getPerimeter() {
        return 2 * (width + height);
    }
};

int main() {
    // Create a rectangle with width 5 and height 3
    Rectangle rect(5.0, 3.0);

    // Output area and perimeter
    cout << "Area: " << rect.getArea() << endl;
    cout << "Perimeter: " << rect.getPerimeter() << endl;

    return 0;
}
\end{lstlisting}

\subsubsection{Inline Code}\label{subsubsec:inline-code}
You can also include inline code with \lstinline|print("Hello, world!")| that fits within your text.

\subsubsection{Configure Listings Styles}\label{subsubsec:listing-styles}
In your preamble, you can define custom listing styles:

\begin{verbatim}
\lstdefinestyle{mystyle}{
    backgroundcolor=\color{lightgray!30},
    basicstyle=\ttfamily\small,
    breakatwhitespace=false,
    breaklines=true,
    captionpos=b,
    keepspaces=true,
    numbers=left,
    numbersep=5pt,
    showspaces=false,
    showstringspaces=false,
    showtabs=false,
    tabsize=2
}
\lstset{style=mystyle}
\end{verbatim}

\subsection{Glossary Usage}\label{subsec:glossary-usage}
LaTeX provides excellent support for creating glossaries through the \texttt{glossaries} package.

\subsubsection{Defining Terms}\label{subsubsec:defining-terms}
In your preamble or in a separate file, define glossary terms:

\begin{verbatim}
\newglossaryentry{latex}{
    name=LaTeX,
    description={A document preparation system}
}

\newglossaryentry{pdf}{
    name=PDF,
    description={Portable Document Format}
}
\end{verbatim}

\subsubsection{Using Glossary Terms}\label{subsubsec:using-terms}
Reference glossary terms in your document with:

\begin{verbatim}
\gls{latex} is used to create \gls{pdf} files.
\end{verbatim}

Different forms are available:
\begin{itemize}
    \item \texttt{\textbackslash gls\{term\}} - Standard reference
    \item \texttt{\textbackslash Gls\{term\}} - Capitalized first letter
    \item \texttt{\textbackslash glspl\{term\}} - Plural form
    \item \texttt{\textbackslash glsadd\{term\}} - Add to glossary without displaying
\end{itemize}

\subsubsection{Printing the Glossary}\label{subsubsec:printing-glossary}
To include the glossary in your document:

\begin{verbatim}
\printglossaries
\end{verbatim}

Remember that you need to run additional commands to generate the glossary:
\begin{verbatim}
pdflatex document
makeglossaries document
pdflatex document
\end{verbatim}

\subsection{Bibliography and Citations}\label{subsec:bibliography}
LaTeX offers powerful bibliography management through BibTeX and BibLaTeX.

\subsubsection{BibTeX Format}\label{subsubsec:bibtex-format}
Bibliography entries are stored in a \texttt{.bib} file:

\begin{verbatim}
@book{knuth1986texbook,
  title={The TeXbook},
  author={Knuth, Donald E},
  year={1986},
  publisher={Addison-Wesley}
}

@article{lamport1994latex,
  title={LaTeX: A Document Preparation System},
  author={Lamport, Leslie},
  journal={Addison-Wesley},
  year={1994}
}
\end{verbatim}

\subsubsection{Citing Sources}\label{subsubsec:citing}
Citations can be added to your document using various commands:

\begin{verbatim}
\cite{knuth1986texbook}     % Standard citation: [1]
\textcite{lamport1994latex} % Author citation: Lamport [2]
\parencite{knuth1986texbook} % Parenthetical: (Knuth, 1986)
\end{verbatim}

\subsubsection{Bibliography Styles}\label{subsubsec:bibliography-styles}
Different citation styles can be used:

\begin{verbatim}
\bibliographystyle{plain}    % Standard numbered
\bibliographystyle{unsrt}    % Unsorted numbered
\bibliographystyle{alpha}    % Alpha-style [Knu86]
\bibliographystyle{ieeetr}   % IEEE Transactions style
\end{verbatim}

\subsubsection{Printing the Bibliography}\label{subsubsec:printing-bibliography}
To include the bibliography in your document:

\begin{verbatim}
\bibliography{references}  % references.bib is your database
\end{verbatim}

For BibLaTeX, the commands are different:

\begin{verbatim}
\printbibliography[title=References]
\end{verbatim}

\subsection{Cross-Referencing}\label{subsec:cross-referencing}
LaTeX makes it easy to reference different parts of your document:

\begin{itemize}
    \item Sections: \texttt{\textbackslash ref\{sec:technical-features\}}
    \item Equations: \texttt{\textbackslash eqref\{eq:einstein\}}
    \item Figures: \texttt{\textbackslash ref\{fig:example\}}
    \item Tables: \texttt{\textbackslash ref\{tab:sample\}}
\end{itemize}

For example, refer to Section~\ref{sec:technical-features} or see Figure~\ref{fig:simple-example} in the visuals section.

\subsection{Hyperlinks}\label{subsec:hyperlinks}
The \texttt{hyperref} package allows adding hyperlinks:

\begin{verbatim}
\href{https://www.latex-project.org/}{LaTeX Project Website}
\end{verbatim}

Example: \href{https://www.latex-project.org/}{LaTeX Project Website}

You can also create internal links with:

\begin{verbatim}
\hyperref[sec:technical-features]{Technical Features Section}
\end{verbatim}